\documentclass[12pt]{article}

\usepackage[margin=1in]{geometry} % 1-inch margins
\usepackage{times}                % Times New Roman font
\usepackage{fancyhdr}             % Custom header/footer
\usepackage{titlesec}             % Custom section formatting
\usepackage{setspace}             % Line spacing
\usepackage{indentfirst}      
\usepackage{graphicx}
\usepackage{caption}

% Indent first paragraph

\pagestyle{fancy}
\fancyhf{} % Clear default header and footer

% Header (Right-aligned)
\fancyhead[R]{\small Saksham AI: Multilingual Virtual Interviewer}

% Footer (Left-aligned: College Name, Right-aligned: Page Number)
\fancyfoot[L]{\small Pillai HOC College of Engineering and Technology (Autonomous), Rasayani}
\fancyfoot[R]{\small \thepage}

% Horizontal lines in header & footer
\renewcommand{\headrulewidth}{0.4pt} 
\renewcommand{\footrulewidth}{0.4pt} 

% Page numbering starts from 1
\setcounter{page}{1}

% Section Formatting
\titleformat{\section}{\Large\bfseries}{}{0em}{} % Main section bold
\titleformat{\subsection}{\large\bfseries}{}{0em}{} % Subsections bold & larger
\titleformat{\subsubsection}{\normalsize}{}{1em}{} % Sub-subsections normal

% Paragraph Formatting
\setlength{\parindent}{1.5em} % Indent first line
\setlength{\parskip}{0em}     % No extra space between paragraphs
\onehalfspacing               % 1.5 line spacing for readability

\begin{document}

% CHAPTER TITLE PAGE
\vspace*{\fill}
\begin{center}
    {\Large \bfseries Chapter 1} \\[2em]
    {\Large \bfseries Introduction}
\end{center}
\vspace*{\fill}
\newpage % Start content on a new page

% INTRODUCTION SECTION
%\begin{center}
%    {\bfseries Chapter 1} \\
%    {\bfseries Introduction}
%\end{center}

\subsection{1.1 Background}

\vspace{10pt}

  \noindent In recent years, the demand for AI-powered and interactive interview training systems has risen significantly with the advancement of virtual communication and remote recruitment technologies. Traditional interview preparation methods often lack personalized feedback, real-time interaction, and adaptive evaluation—key elements necessary for developing communication and professional readiness skills. Consequently, AI-driven education and assessment systems have emerged as powerful solutions, allowing learners to experience realistic interview simulations that combine theory, practice, and evaluation within a single digital environment. Many existing interview preparation tools focus only on static question banks or language training, leaving a gap between conceptual readiness and real-world interview performance. Similarly, evaluation and feedback processes are often managed on separate systems, leading to fragmented and inefficient learning experiences. To address these challenges, there is a growing need for a Saksham AI-multilingual virtual interviewer platform that supports real-time interaction, speech analysis, automated feedback, guided learning, and performance-based assessments — all under one system.
  
\subsection{1.2 Relevance}
\vspace{10pt}

   \noindent We are driven by the belief that interview preparation should be accessible, intelligent, and inclusive for candidates from all linguistic and educational backgrounds. Our platform seeks to eliminate the barriers that hinder confident communication during interviews. By providing a multilingual AI-powered interviewer and real-time response evaluation, we aim to empower candidates to practice effectively, improve their communication skills, and gain confidence in facing real interview situations. Moreover, we recognize that the recruitment industry is rapidly adopting Artificial Intelligence to improve efficiency and enhance candidate assessment, and we are determined to contribute to this transformation. The increasing competition for jobs and the growing importance of communication and soft skills have further highlighted the need for innovative interview preparation solutions. With Saksham AI – Multilingual Virtual Interviewer, we aim to bring this vision to life by creating an interactive platform that provides realistic interview simulations, automated feedback, and performance analysis. Our goal is to set new standards in AI-driven interview training and ensure that every candidate’s preparation journey is smooth, engaging, and skill-focused. 
   \vspace{5pt}
   
\subsection{1.3 Organization of the Report}
\vspace{5pt}
\noindent This report is organized into several chapters, each explaining different aspects of the Saksham AI: Multilingual Virtual Interviewer project in a structured and systematic manner: \\[2pt]

\noindent Chapter 1: Introduction

\noindent This chapter gives an overview of the Saksham AI: Multilingual Virtual Interviewer, covering its background, relevance, and purpose. It highlights the system’s role in improving virtual communication, interview readiness, and global employability through AI-driven interaction.
 \\[2pt]

\noindent Chapter 2: Literature Review

\noindent  This section reviews related AI interviewing and language processing systems, identifying gaps like limited personalization, lack of real-time feedback, and poor multilingual support that the proposed system aims to overcome.
 \\[2pt]

\noindent Chapter 3: Planning and Requirement Gathering

\noindent  This chapter outlines the project plan, requirements, and technologies used, including NLP, speech recognition, and emotion analysis, which power the system’s intelligent and multilingual interaction features.
 \\[2pt]

\noindent Chapter 4: Project Analysis

\noindent  This chapter analyzes the system through problem identification, feasibility study, and use case analysis, explaining the workflow, user roles (interviewee, admin, evaluator), and key challenges in multilingual processing and performance evaluation.  \\[2pt]   

\noindent Chapter 5: Project Design

\noindent  This section outlines the system use case diagram. It shows how modules like AI Interview Engine, Speech-to-Text, Translation, and Evaluation interact seamlessly within the platform. \\[2pt]

\noindent Chapter 6: Proposed System

\noindent  This chapter describes the main functions and features of the Multilingual Virtual Interviewer, highlighting real-time analysis, multilingual support, emotional feedback, and secure proctoring that enhance the overall interview training experience. \\[2pt]
	
\noindent Chapter 7: Implemented System

\noindent  This chapter explains the development and implementation of the system, including the design of the user interface, integration of multilingual interaction, AI-based question generation, response evaluation, and database management that enable smooth functioning of the multilingual AI interviewer platform. \\[2pt]

\noindent Chapter 8: Result Analysis

\noindent  This chapter presents the evaluation and analysis of the system’s performance, highlighting user response accuracy, feedback effectiveness, language adaptability, and overall system efficiency that demonstrate the reliability and usefulness of the AI-based interview preparation platform.  \\[2pt]

\noindent Chapter 9: Conclusion and Future Scope

\noindent  This chapter summarizes the overall achievements of the project and discusses possible future enhancements, including improved AI intelligence, advanced speech analysis, expanded language support, and additional training features that can further strengthen the effectiveness of the multilingual AI interviewer system.  \\[2pt]


\end{document}